\tinysidebar{\debug{exercises/{37-08-if-part2/question.tex}}}Get an integer n from the user and then compute the

sum of the rightmost 4 digits of n. Once a digit 0 is added, you should

stop the summation. Do this with repetition code. Here's a test case



\begin{consolethree}[escapeinside=||]

|\textbf{1534315}|



13 

\end{consolethree}





Here's another:



\begin{consolethree}[escapeinside=||]

|\textbf{1534015}|



6 

\end{consolethree}





(after adding \texttt{5} and \texttt{1}, on seeing a \texttt{0}, the summation

stops. Therefore the sum is 6) and another





\begin{consolethree}[escapeinside=||]

|\textbf{1534205}|



5 

\end{consolethree}



Here's the idea. First here's the code to sum the 4 rightmost digits:



\begin{console}

#include <iostream>



int main()

{

    int n;

    std::cin >> n;



    int s = 0;



    s = s + n % 10;



    n = n / 10;



    s = s + n % 10;



    n = n / 10;



    s = s + n % 10;



    n = n / 10;



    s = s + n % 10;



    n = n / 10;



    std::cout << s << n;



    return 0;

} 





Each block needs to know if a zero was encountered. So \ldots



\begin{console}

#include <iostream>



int main()

{

    int n;

    std::cin >> n;



    int s = 0;

    bool zero_not_found = true;



    if (zero_not_found)

    {

        s = s + n % 10;

        n = n / 10;

    }



    if (zero_not_found)

    {

        s = s + n % 10;

        n = n / 10;

    }



    if (zero_not_found)

    {

        s = s + n % 10;

        n = n / 10;

    }



    if (zero_not_found)

    {

        s = s + n % 10;

        n = n / 10;

    }



    std::cout << s << ", " <<n;

    return 0;

} 





Get it? Think about what will happen if you don't use this communication device \verb!zero_not_found!. You will most likely have a deeply nested if-else structure in your code:



\begin{console}

if (n % 10 != 0) // 1st digit != 0

{

    s = s + n % 10;

    n = n / 10;

    if (n % 10 != 0) // 2nd digit != 0

    {

        s = s + n % 10;

        n = n / 10;

        if (n % 10 != 0) // 3rd digit != 0

        {

            s = s + n % 10;

            n = n / 10;

            if (n % 10 != 0) // 4th digit != 0

            {

                s = s + n % 10;

                n = n / 10;

            }

        }

    }

}

\end{console}

